\section{Integrated Discussion}\label{sec:discussion}

We asked to what extent SAEs preserve the information and behaviour of a language model beyond reconstructing its
activations well. The answer from this study is that reconstruction quality is necessary evidence of faithfulness but
is insufficient by itself: it tracks behavioural faithfulness reasonably well, it tracks information and causal
preservation poorly, and improvements at one level do not reliably transfer to the others.

\subsection{What reconstruction quality does tell us}\label{sec:disc-does}

Reconstruction quality is informative. Within a layer of GPT-2, SAEs with lower reconstruction error damage next-token
prediction less, and this association is robust to architecture and survives correction for multiple testing
(Section~\ref{sec:behaviour}). Across depth, reconstruction and splice damage degrade together. The improvement
experiments point the same way: retraining with about a thousand times more data and four times the width cut
reconstruction error by 31--74\,\% and splice damage by 35--79\,\% in a pre-registered held-out test
(Section~\ref{sec:r2}). Within comparable settings, then, better reconstruction is a good sign for behavioural
faithfulness, and a large reconstruction improvement is likely to bring a behavioural improvement with it.

\subsection{What reconstruction quality does not tell us}\label{sec:disc-doesnot}

Reconstruction quality does not guarantee information preservation. Within a layer it does not robustly predict
truth-information loss, persona information is lost by SAEs that reconstruct well, and the large reconstruction gains
of R2 brought no confirmed gain in truth information (Sections~\ref{sec:info} and~\ref{sec:r2}). It does not
guarantee causal preservation either: at the later GPT-2 sites the reconstruction carries about two-thirds of a causal
effect the model uses, and at layer 11 the architecture with higher explained variance retains less of that effect
(Section~\ref{sec:causal}). Nor does it guarantee robustness under targeted optimisation: the behaviour-targeted R4
fine-tune lowered splice damage while slightly worsening reconstruction and reducing causal preservation at one site
(Section~\ref{sec:r4}). The common explanation is that reconstruction metrics weight directions by their variance,
whereas task information and causal effects often live in small differences between nearly identical activations. An
SAE can reconstruct the honest and the dishonest prompt of a pair well individually and still reconstruct their
difference poorly, because the two reconstruction errors do not cancel.

\subsection{Why reader choice matters}\label{sec:disc-readers}

Conclusions about information preservation depend on the reader. A probe fitted on original activations and applied
frozen to reconstructions mixes two questions: whether the information is still present, and whether the reconstruction
stays where that probe's decision boundary was fitted. On truth statements the second question dominates, since
norm-matched noise reproduces most of the frozen-probe collapse. The same reader also produced the only apparent truth
gain of fine-tuning, which re-trained readers did not reproduce (Section~\ref{sec:finetune}). Re-trained readers ask
the stricter question of whether the information is still recoverable, and they show a smaller but real truth loss and
a large persona loss. The comparison with PCA adds a second dependence: matched on reconstruction error, PCA keeps more
truth information than the SAEs, whereas matched on the per-token active budget, no trained SAE is classified worse and
some are better. Information-preservation claims should therefore name the reader and the constraint under which a
reference is matched.

\subsection{Faithfulness is not one scalar}\label{sec:disc-scalar}

The improvement experiments make the multi-level structure most visible, because each method improved what it
optimised and little else. Truth-focused fine-tuning lowered task reconstruction error and the frozen-probe gap, but not
the re-trained gap or causal preservation. Data and width scaling improved reconstruction and behaviour on held-out
data, and causal preservation on validation data, but not truth information. A behavioural objective at fixed sparsity
improved behaviour at a small cost in reconstruction and some cost in causal preservation. Where an output-KL term
appeared to improve causal preservation in the earlier fine-tunes, the gain coincided with extra active latents and
vanished when sparsity was fixed. These dissociations are the main evidence that faithfulness is multi-dimensional. A
claim that one SAE is ``more faithful'' than another is only meaningful when it names the level, and an evaluation that
reports only reconstruction and sparsity leaves three of the four levels unmeasured.

\subsection{Cross-model generalisation}\label{sec:disc-gemma}

The Gemma experiment strengthens the behavioural part of this account. In a larger model, with an independently
trained public SAE family, a different hook convention and a different corpus, SAE reconstruction error was again more
harmful to next-token prediction than random error of the same norm, at all three pre-registered layers
(Section~\ref{sec:gemma}). The finding that reconstruction error is not behaviourally neutral is therefore not specific
to GPT-2 or to our training recipe. Gemma does not, however, establish that the information or causal findings
generalise, because it was not evaluated at those levels, and the differences between the two settings rule out any
statement that one model's SAEs are more faithful than the other's.
